%% kapa: exam=1 round=31 year=2020 period=2 form=A subject=accounting subject_ko=회계학 seq=4 number=44
%% source: Q-Net 감정평가사 기출문제 (cst00309) artlSeq=5207219 file=extracted/1008/5207219/A형/A형/2020년도 제31회 감정평가사 1차 시험_2교시_A형.pdf
%% conversion: kordoc 4.24.51 -> tools/gen1.py (auto); review=proofread-sonnet; flags=none
\begin{question}{44}
(주)감평은 20x1년 4월 1일에 만기가 20x1년 7월 31일인 액면금액 ￦1,200,000의 어음을 거래처로부터 수취하였다. (주)감평은 동 어음을 20x1년 6월 30일 은행에서 할인하였으며, 할인율은 연 12\%이다. 동 어음이 무이자부인 어음일 경우(A)와 연 9\%의 이자부어음일 경우(B) 각각에 대해 어음할인시 (주)감평이 금융상품(받을어음)처분손실로 인식할 금액은? (단, 어음할인은 금융상품의 제거요건을 충족 시킨다고 가정하며, 이자는 월할계산한다.)
{\par\noindent\hspace*{2em}\makebox[3cm][c]{\uline{(A)}}\makebox[3cm][c]{\uline{(B)}}\par}
\begin{choices}
  \item \makebox[3cm][c]{￦0}\makebox[3cm][c]{￦3,360}
  \item \makebox[3cm][c]{￦0}\makebox[3cm][c]{￦12,000}
  \item \makebox[3cm][c]{￦12,000}\makebox[3cm][c]{￦3,360}
  \item \makebox[3cm][c]{￦12,000}\makebox[3cm][c]{￦9,000}
  \item \makebox[3cm][c]{￦12,000}\makebox[3cm][c]{￦12,000}
\end{choices}
\end{question}
\answer{3}
